% Page 069

\seite{69}

\margmark{wertbestän\\Mark hat n\\1.000.000.\\Löhne und\\geb. vor e\\rechnung b\\Zu dem Gol\\extra 4 Pf\\Die Briefe}%
diges Zahlungsmittel, die Rentenmark. Die deutsche\ob
un aufgehört, durch vorschr. Billionen Mark =\ob
000.000 d.\ 1 Renten Mark oder 1 Gold Mark.\ob
Gehälter werden nach Goldmark berechnet, vorgestern\ob
inigen Monaten einen Festpreis in der Goldbe\ohb
estätigt. Die Zahlungskraft des Volks ist erschüttert.\ob
dmark-Briefpost wurde festg. dem 5 Pf.\ w. 1 Monta\unsicher\ob
g. U. der Preis derselben Dinge geht vor der Mark.\ob
s Postpost ist als zu verwöhnlich.

Hoffen wir, daß das neue Jahr uns ruhig eine bessere\ob
Lage bringe, daß wir wieder mit Freund zu allen Dingen\ob
und wahrem Wissen schaffen können!

\margmark{20. Januar}%
.   Die Ausgabe für einen starken Winter haben nicht ge\ohb
bracht. Manche Rübenvorräte drohten den Ochsen, alles\ob
zugefroren, die gestörten Rüben können von den Leuten\ob
nicht d{[em]} Vieh verfüttert werden. Ein strenger Winter!

\margmark{22. Januar}%
.   In der gestrigen Gemeindeversamml. und Gemeinderechnungsprüfung\ob
wurde beschlossen, daß die Einkommensteuer und Steuern\ob
in der Hälfte verteilen vorgesetzt und nicht mehr durch\ob
die Ausgabe ausgeschickt werden soll.

\pend
\begin{verse}
Die Hebesonnierung wird um zweiter Hälfte beseitigt.
\end{verse}
\pstart

\margmark{2. Februar}%
.   Der Gemeind. hat die gestellt Ober für die Befsondungng\unsicher genehmigt.\ob
Am 2.\ Februar erhielt die Hälfte ihren Antrag durch die Gemde Bürger\ohb
meister und Herbert Kley und Michael Mattscheck\unsicher. Der Beschluß der Hälfte\ob
Steuer geht ohne Festbelassung. Die Kapitalertragssteuer wird 150%-\ob
ige und die gleiche ist zum 3ten für alle auszusprechen.

\pend
\begin{verse}
Die Gemeinde sollte denn schuldigsten und ärmsten Oaster
\end{verse}
\pstart
Gehalts der Gemeindeordnungsmäßigen Beschlüßf ane Vorteil ohne\ob
Nachteile gesprochen wird sich selbst abschaffen wobei man nur\ob
höchste der Rechnungl\unsicher behalten, die so, am bk. Verwendungm\unsicher ist, und\ob
zum Teil nicht der seinen anderen!

\pend
\begin{verse}
Der Pfarrer scheint sich nicht mehr zu wollen. Tägliche
\end{verse}
\pstart
Verkehrung man sich fiel und Unfreundl. ein Boden und fal\ohb
sches still genuggenügen macht Eltern in großer Menge, der Ver\ohb
br. Sache – Lust die Ältern die Lust, für die Jugend zu sorgen.\ob
Ein glückliches Himmel scheint zu befestigen zu werden. Die Grenze\ob
Versuchungen ist schrecklich, die Plane mußallen\unsicher{} dagegen, die in Ent\unsicher
\pend
